% REPORT_TEMPLATE.tex — CA22 LaTeX report template
\documentclass[11pt]{article}
\usepackage[utf8]{inputenc}
\usepackage{geometry}
\geometry{margin=1in}
\usepackage{amsmath,amssymb}
\usepackage{graphicx}
\usepackage{booktabs}
\usepackage{caption}
\usepackage[round]{natbib}
\usepackage{hyperref}
\hypersetup{colorlinks=true,linkcolor=blue,citecolor=blue,urlcolor=blue}

% Short macros for convenience
\newcommand{\todo}[1]{{\color{red} TODO: #1}}

\title{\textbf{Title of Your Experiment}\vspace{0.5em}\\\large CA22 — Reproducible Experiment Template}
\author{Author Name(s) \\\texttt{email@example.com}}
\date{\today}

\begin{document}
\maketitle

\begin{abstract}
A concise abstract (approx. 150–300 words): motivation, method, and key quantitative results. Keep this short and self-contained.
\end{abstract}

\section{Introduction}
Provide a short, precise problem statement and motivation. Cite relevant background briefly (e.g., \citep{Sutton2018}).

\section{Related Work}
One or two paragraphs placing your experiment in context with the most relevant prior work.

\section{Methods}
Describe the model and objectives. Include pseudocode or equations for losses. Example subsections:
\subsection{Model}
Describe the policy and value architectures you used (layer sizes, activations) and reference `src/model.py`.

\subsection{Objective and Constraints}
Describe the loss functions, including any Lagrangian formulation (reference `src/losses.py`) and the update rules used for Lagrange multipliers.

\subsection{Dataset and Synthetic Generator}
Describe the synthetic generator (`src/data.py`), episode length, and any dataset sampling details.

\section{Experimental Setup}
List hyperparameters and training details (optimizer, learning rate, batch size, epochs). Add a small table if useful.

\section{Results}
Present your quantitative results with figures and tables. Include clear captions and short interpretation.

\begin{figure}[ht]
  \centering
  % include a PNG or PDF plot here: outputs/figures/example.png
  \fbox{\parbox[c][4cm][c]{0.9\linewidth}{\centering Placeholder for a figure (save as outputs/figures/your_figure.png)}}
  \caption{Example figure caption explaining what is shown and the main takeaway.}
\end{figure}

\section{Discussion}
Interpret results and state limitations, potential failure modes, and suggestions for future work.

\section{Reproducibility Appendix}
List exact commands and configuration used to reproduce the experiment. Example:
\begin{verbatim}
python -m venv .venv && source .venv/bin/activate
pip install -r requirements.txt
python run_experiment.py --config configs/my_experiment.yaml
\end{verbatim}
Include a table of key hyperparameters and seeds.

\section*{Acknowledgements}
(Optional) Acknowledge assistance or funding.

\bibliographystyle{plainnat}
\bibliography{refs}

\end{document}
